We study whether ETF-linked information can improve daily stock return forecasting in U.S. equities when performance is judged in economic rather than purely statistical terms. The empirical design combines rolling-window estimation, Ridge-based return prediction, and an explicit execution layer with holding-period and rebalance controls. The starting point is negative out-of-sample $R^2$: large raw ETF feature blocks perform poorly and do not reliably outperform a stock-only HAR baseline. The evidence instead indicates that ETF-linked information is most useful when it enters through small, stock-specific adjustments rather than as a large covariate block. On that basis, we construct three strategies: a consensus-majority ETF signal, a regime-dependent ETF refinement, and a conservative baseline-plus-top3-ETF blend. We evaluate them with AlphaMark and with separate transaction-cost sensitivity checks at 10, 20, and 50 basis points. The best gross AlphaMark slices reach Sharpe ratios around 0.5--0.6, and the strongest specifications remain competitive with an optimized baseline after cost adjustments. Overall, ETF information is not effective as a raw feature expansion, but it can become economically useful when incorporated through low-dimensional signal design and disciplined execution.
